
This work synthesizes insights from software package deprecation, dataset documentation standards, and machine learning repository design to address gaps that none of these domains solve in isolation.

\subsubsection{Software Package Deprecation}

Package managers (NPM, PyPI, Maven) provide automated deprecation mechanisms via version constraints and load-time warnings. NPM's deprecate command marks packages obsolete and displays warnings during installation. Python's DeprecationWarning system flags obsolete APIs at import time. These systems demonstrate that automated, point-of-use deprecation notices are feasible with minimal overhead.

Software package deprecation relies on semantic versioning (v1.0 $\rightarrow$ v2.0) and assumes single-path linear succession: one canonical replacement per deprecated package. This model does not capture task-conditional successors in machine learning datasets. ImageNet's successor varies by use case (ImageNet-v2 for robustness testing versus ImageNet-21k for pretraining). The proposed system extends package manager patterns with context-aware successor graphs modeling multi-path, task-specific replacement relationships.

\subsubsection{Dataset Documentation and Lifecycle Management}

Datasheets for Datasets \citep{Gebru2018} and Data Statements \citep{bender2018data} establish documentation standards for machine learning datasets, covering motivation, composition, collection process, and maintenance procedures. These frameworks improve transparency but operate as static documentation rather than executable infrastructure. Datasheets do not detect deprecation candidates automatically, recommend successors at load time, or measure adoption outcomes.

The proposed instrumented policies complement documentation approaches by adding automated enforcement: health metrics surface deprecation candidates without manual curator review, context-aware graphs provide personalized successor recommendations, and load-time telemetry enables quantitative adoption measurement.

\subsubsection{ML Repository Design and Versioning}

HuggingFace Datasets Hub, OpenML, and UCI Machine Learning Repository support dataset sharing but lack formal deprecation infrastructure. HuggingFace uses Git-based versioning with manual README updates for deprecation notices. OpenML relies on dataset status flags and maintainer annotations. UCI provides static archives with no version tracking. All three platforms depend on documentation-only signals that users must actively discover.

Recent work on FAIR principles for machine learning datasets emphasizes findability, accessibility, interoperability, and reusability but does not address lifecycle management post-publication. Versioning exists (Git tags, dataset revisions) but lacks deprecation-specific mechanisms: no automated health metrics, no task-aware successor recommendations, and no adoption tracking.

The proposed system builds on HuggingFace's programmatic loader infrastructure to add three missing components validated through proof-of-concept experiments. Health metrics adapt software's semantic versioning to machine learning's usage-driven lifecycle (velocity, emergence, issue signals replace version numbers). Context-aware graphs capture task-specific needs that package managers miss. Load-time instrumentation provides measurement infrastructure that documentation alone cannot.

\subsubsection{Positioning}

No prior system combines automated detection, context-aware recommendation, and adoption tracking for machine learning dataset deprecation. Software package managers provide the closest template but lack task-specific successor modeling. Dataset documentation improves transparency but not automation. Machine learning repositories enable sharing but not lifecycle management. This work fills the infrastructure gap through validated mechanistic integration of all three components.
